\section{Conclusion}
\label{sec:conclusion}

We have introduced the Moving-Knife Algorithm (\mka{}), the first redistricting
structure algorithm that explicitly maximises Reock compactness.
\mka{} adapts the fair-division moving-knife protocol of Dubins and
Spanier~\citeyear{dubins1961} to the bisection framework, sweeping through 180
orientations and selecting the one that maximises the minimum Reock score of
the two population-balanced halves.
Implementation reuses the Albers EPSG:5070 projection and centroid companion
file from \cvdgeo{} (T.11~\citep{dellaLibera2026cvdgeo}), and the 200-iteration
boundary-swap rebalance shared with all geometric structure algorithms.

\paragraph{When to use \mka{}.}
Use \texttt{--structure moving-knife} when:
\begin{itemize}
\item The jurisdiction's constitutional or statutory standard explicitly
  cites Reock compactness (North Carolina, Florida, Wisconsin --- see
  Section~\ref{sec:introduction}).
\item A court order or remedial mandate specifies that the submitted plan
  must demonstrate Reock maximisation.
\item A procedurally neutral, legally articulable justification is needed
  that is more accessible to non-technical judges than the graph-theoretic
  arguments of \metis{} or the Lloyd-iteration arguments of \cvdgeo{}.
\end{itemize}

\paragraph{When other algorithms may be preferable.}
\begin{itemize}
\item For Polsby-Popper optimisation: use \cvdgeo{}
  (\texttt{--structure centroidal-voronoi --cvd-metric geographic}).
\item For minimum edge cut: use \metis{}
  (\texttt{--structure standard-bisect}).
\item For geographic balance speed: use \bfsg{}
  (\texttt{--structure bfs-growth}).
\item When Reock is not the operative metric: \mka{}'s 10--15\,\% EC premium
  over \cvdgeo{} may not be justified.
\end{itemize}

\paragraph{Position in the algorithm portfolio.}
\mka{} joins \cvdgeo{} (T.11) and \bfsg{} (T.12~\citep{dellaLibera2026bfs})
as the third geometric structure algorithm in the platform.
Together, the three algorithms provide coverage of the three compactness
metrics most commonly cited in redistricting litigation:
\bfsg{} targets geographic balance (used as a proxy for Convex Hull Ratio),
\cvdgeo{} targets Polsby-Popper via geographic proximity, and
\mka{} targets Reock via angular sweep.
Practitioners selecting among them should match the algorithm to the legally
operative metric in their jurisdiction.

\paragraph{Open questions.}
Several extensions are deferred:
\begin{enumerate}
\item \textbf{Polygon-boundary MEC (Phase~2)}: replace the centroid
  approximation (Remark~\ref{rem:centroid-approx}) with the full
  polygon-boundary MEC using the TIGER/Line shapefile vertex data.
  This will eliminate the $\reock{} > 1$ clamping cases and produce
  exact Reock scores.

\item \textbf{$n_{\mathrm{orient}}$ tradeoff}: benchmark 36 vs.\ 180 orientations
  on FL and WA panhandle geometries to determine whether the 5$\times$
  speedup of 36 orientations is worth the small quality loss.

\item \textbf{\mka{} + BisectionEnsemble}: run \be{}$(T=50)$ with \mka{}
  bisectors to find the best Reock-maximising seed among 50 independent
  runs.
  Per-run runtime ($\approx$6\,s on NC) makes the ensemble feasible with
  \texttt{--workers 8}.

\item \textbf{Non-convex subgraphs}: for states with severe non-convexity
  (peninsulas, island tracts), investigate whether a post-hoc local search
  over orientations --- seeded by \mka{}'s $\theta^*$ --- can improve
  on the linear sweep's heuristic.
\end{enumerate}

\paragraph{Summary.}
The Moving-Knife Algorithm provides a principled, court-defensible
procedure for Reock maximisation in congressional redistricting.
Its theoretical basis in fair-division theory (Dubins-Spanier 1961,
Puppe-Tasnadi 2026) provides academic grounding; its empirical performance
(highest Min-Reock on NC, FL, and WA) confirms the operational claim;
and its fair-division certificate offers a procedurally neutral justification
that is accessible to non-technical decision-makers.
The algorithm is available in the \texttt{bisect} platform as
\texttt{--structure moving-knife}.
